%% Online quizzes, AP Physics 1 Unit 2 (Force and Translational Dynamics), aligned with the guided notes 2.1 to 2.9.
\renewcommand{\nocalc}{}
%% Blocks: one per topic (Q2.1 ... Q2.9) and the unit test (QTEST). Use g = 10 m/s^2.
\newcommand{\ms}{\,\mathrm{m/s}}
\newcommand{\mss}{\,\mathrm{m/s^2}}
\newcommand{\m}{\,\mathrm{m}}
\newcommand{\s}{\,\mathrm{s}}
\newcommand{\N}{\,\mathrm{N}}
\newcommand{\kg}{\,\mathrm{kg}}
\newcommand{\Nm}{\,\mathrm{N/m}}

%%DOC Q2.1 | Systems and Center of Mass
\begin{exo}{}
Two $2\kg$ balls are at $x=0$ and $x=6\m$. Where is their center of mass?
\choices{$x=3\m$}{$x=6\m$}{$x=2\m$}{$x=12\m$}
\answer{A}{Equal masses: halfway.}
\end{exo}
\begin{exo}{}
A $1\kg$ mass is at $x=0$ and a $2\kg$ mass is at $x=3\m$. Where is their center of mass?
\choices{$x=1.5\m$}{$x=2\m$}{$x=1\m$}{$x=3\m$}
\answer{B}{$\frac{1(0)+2(3)}{3}=2\m$.}
\end{exo}
\begin{exo}{}
Internal forces in a system
\choices{are always zero}{are external forces}{cannot change the motion of the center of mass}{always speed the system up}
\answer{C}{They come in equal and opposite pairs inside the system.}
\end{exo}
\begin{exo}{}
A projectile breaks into two pieces in mid-air (no air resistance). The path of the center of mass of the two pieces is
\choices{ended}{the same parabola as before}{straight down}{a horizontal line}
\answer{B}{Only gravity acts on the system, before and after.}
\end{exo}
\begin{exo}{}
When can a runner be modeled as a single point object?
\choices{Never}{Only when at rest}{Only when spinning}{When only the motion of the whole body matters}
\answer{D}{If arms and legs do not matter for the question, the object model works.}
\end{exo}
\begin{exo}{}
The system is ``a book and a table''. Which force is internal?
\choices{Earth's gravity on the book}{The floor pushing on the table}{The book pushing on the table}{Earth's gravity on the table}
\answer{C}{Both the book and the table are inside the system.}
\end{exo}

%%DOC Q2.2 | Forces and Free-Body Diagrams
\begin{exo}{}
What is the weight of a $0.5\kg$ apple?
\choices{$0.5\N$}{$5\N$}{$50\N$}{$0.05\N$}
\answer{B}{$mg=0.5(10)=5\N$.}
\end{exo}
\begin{exo}{}
The normal force on an object is
\choices{always vertical and upward}{parallel to the surface}{perpendicular to the surface, pushing away from it}{directed toward the center of Earth}
\answer{C}{``Normal'' means perpendicular.}
\end{exo}
\begin{exo}{}
A tension force is exerted by
\choices{a rope, string or cable}{a spring}{a surface}{Earth}
\answer{A}{Tension pulls along a rope or string.}
\end{exo}
\begin{exo}{}
A book slides to the right on a table and slows down. Which free-body diagram is correct?
\choices{Weight, normal force, friction to the left}{Weight, normal force, a ``force of motion'' to the right}{Weight only}{Weight, normal force, friction to the right}
\answer{A}{No force pushes it forward; friction opposes the sliding.}
\end{exo}
\begin{exo}{}
An object falls downward through the air. In which direction is the air resistance?
\choices{Down}{Up}{Sideways}{There is none}
\answer{B}{Air resistance opposes the velocity.}
\end{exo}
\begin{exo}{}
A skydiver falls at a constant speed. Compare the air resistance with the weight.
\choices{It is zero.}{It is smaller.}{It is larger.}{It is equal.}
\answer{D}{Constant velocity: the net force is zero.}
\end{exo}

%%DOC Q2.3 | Newton's Third Law
\begin{exo}{}
A rocket pushes hot gas downward. The gas pushes the rocket
\choices{downward}{upward}{not at all}{sideways}
\answer{B}{Third-law pair: opposite direction.}
\end{exo}
\begin{exo}{}
The two forces of a third-law pair are
\choices{on the same object}{unequal}{equal, opposite, and on different objects}{in the same direction}
\answer{C}{That is why they never cancel.}
\end{exo}
\begin{exo}{}
A mosquito hits the windshield of a truck. Compare the force on the mosquito with the force on the truck.
\choices{Equal in magnitude}{Larger on the mosquito}{Larger on the truck}{Zero on the truck}
\answer{A}{Third law: always equal.}
\end{exo}
\begin{exo}{}
In the same collision, which has the larger acceleration?
\choices{The truck}{The mosquito}{They are equal.}{Neither}
\answer{B}{Same force, much smaller mass.}
\end{exo}
\begin{exo}{}
You push on a wall with $40\N$. The wall pushes on you with
\choices{$0$}{$80\N$}{$20\N$}{$40\N$, toward you}
\answer{D}{Equal and opposite.}
\end{exo}
\begin{exo}{}
The floor pushes up on you (normal force). What is the third-law partner of this force?
\choices{Your weight}{Your push down on the floor}{Earth's pull on you}{The floor pushing on Earth}
\answer{B}{Same two objects (you and floor), same type (contact force), opposite direction.}
\end{exo}

%%DOC Q2.4 | Newton's First Law
\begin{exo}{}
An object at rest has zero net force on it. It will
\choices{start to move}{stay at rest}{accelerate}{fall}
\answer{B}{First law.}
\end{exo}
\begin{exo}{}
A hockey puck slides on frictionless ice after being hit. It
\choices{moves at constant velocity}{slows down}{speeds up}{stops at once}
\answer{A}{No net force, no change in velocity.}
\end{exo}
\begin{exo}{}
The mass of an object measures its
\choices{speed}{volume}{inertia}{weight}
\answer{C}{Mass is the measure of inertia. Weight is a force.}
\end{exo}
\begin{exo}{}
A $20\kg$ crate hangs from a crane cable and moves down at a constant speed. What is the tension?
\choices{$0$}{$200\N$}{Less than $200\N$}{More than $200\N$}
\answer{B}{Constant velocity: $F_T=mg$.}
\end{exo}
\begin{exo}{}
Two perpendicular forces of $5\N$ and $12\N$ act on an object. What third force keeps it in equilibrium?
\choices{$7\N$}{$17\N$}{$13\N$}{$60\N$}
\answer{C}{It must cancel the resultant $\sqrt{5^2+12^2}=13\N$.}
\end{exo}
\begin{exo}{}
A car moves at a constant $20\ms$ on a straight road. The net force on it is
\choices{forward}{backward}{equal to the friction}{zero}
\answer{D}{Constant velocity.}
\end{exo}

%%DOC Q2.5 | Newton's Second Law
\begin{exo}{}
A net force of $10\N$ acts on a $2\kg$ object. What is its acceleration?
\choices{$20\mss$}{$5\mss$}{$0.2\mss$}{$12\mss$}
\answer{B}{$10/2=5\mss$.}
\end{exo}
\begin{exo}{}
What net force gives a $1000\kg$ car an acceleration of $3\mss$?
\choices{$3000\N$}{$333\N$}{$1003\N$}{$30\,000\N$}
\answer{A}{$F=ma=3000\N$.}
\end{exo}
\begin{exo}{}
The net force on an object doubles; its mass is unchanged. Its acceleration
\choices{is halved}{stays the same}{doubles}{quadruples}
\answer{C}{$a$ is proportional to the net force.}
\end{exo}
\begin{exo}{}
A $60\kg$ person rides an elevator accelerating upward at $1\mss$. What is the normal force on her?
\choices{$540\N$}{$600\N$}{$60\N$}{$660\N$}
\answer{D}{$60(10+1)=660\N$.}
\end{exo}
\begin{exo}{}
An elevator's cable breaks and it falls freely. What does a scale in the elevator read?
\choices{$0$}{$mg$}{$2mg$}{$mg/2$}
\answer{A}{In free fall, $F_N=m(g-g)=0$.}
\end{exo}
\begin{exo}{}
A $2\kg$ block and a $3\kg$ block touch on a frictionless floor and are pushed together by $20\N$. What is their acceleration?
\choices{$10\mss$}{$4\mss$}{$6.7\mss$}{$20\mss$}
\answer{B}{System: $20/5=4\mss$.}
\end{exo}

%%DOC Q2.6 | Gravitational Force
\begin{exo}{}
Two objects at distance $r$ attract with force $F$. At distance $2r$ the force is
\choices{$F/2$}{$F/4$}{$2F$}{$4F$}
\answer{B}{Inverse square law.}
\end{exo}
\begin{exo}{}
Both masses are doubled, at the same distance. The force becomes
\choices{$2F$}{$F$}{$4F$}{$F/4$}
\answer{C}{$2\times2=4$.}
\end{exo}
\begin{exo}{}
What is the mass of a $10\kg$ rock on the Moon?
\choices{$10\kg$}{$1.6\kg$}{$16\kg$}{$0$}
\answer{A}{Mass is the same everywhere.}
\end{exo}
\begin{exo}{}
What is the weight of the same rock on Earth?
\choices{$10\N$}{$1000\N$}{$1\N$}{$100\N$}
\answer{D}{$10(10)=100\N$.}
\end{exo}
\begin{exo}{}
What is the gravitational field at a height of $3$ Earth radii above the surface?
\choices{$2.5\,\mathrm{N/kg}$}{$3.3\,\mathrm{N/kg}$}{$0.625\,\mathrm{N/kg}$}{$1.1\,\mathrm{N/kg}$}
\answer{C}{Distance to the center $=4R$: $10/16=0.625\,\mathrm{N/kg}$.}
\end{exo}
\begin{exo}{}
The gravitational field $g$ is measured in
\choices{N}{N/kg (the same as m/s$^2$)}{kg}{J}
\answer{B}{$g=F_g/m$.}
\end{exo}

%%DOC Q2.7 | Kinetic and Static Friction
\begin{exo}{}
A $20\kg$ box slides on a level floor with $\mu_k=0.1$. What is the friction force?
\choices{$2\N$}{$20\N$}{$200\N$}{$0.1\N$}
\answer{B}{$0.1(200)=20\N$.}
\end{exo}
\begin{exo}{}
A box at rest can take up to $40\N$ of static friction. You push it with $15\N$. The friction force is
\choices{$40\N$}{$0$}{$15\N$}{$25\N$}
\answer{C}{Static friction matches the push.}
\end{exo}
\begin{exo}{}
How do $\mu_s$ and $\mu_k$ usually compare?
\choices{$\mu_s\ge\mu_k$}{$\mu_s<\mu_k$}{They are always equal.}{$\mu_k=0$}
\answer{A}{It is harder to start sliding than to keep sliding.}
\end{exo}
\begin{exo}{}
A box sliding on a level floor slows down at $3\mss$ with no push. What is $\mu_k$?
\choices{$3$}{$0.03$}{$0.3$}{$30$}
\answer{C}{$a=\mu_kg\Rightarrow\mu_k=3/10=0.3$.}
\end{exo}
\begin{exo}{}
A sliding $10\kg$ box ($\mu_k=0.2$) is pushed with $50\N$. What is its acceleration?
\choices{$5\mss$}{$3\mss$}{$7\mss$}{$2\mss$}
\answer{B}{$f_k=20\N$; $(50-20)/10=3\mss$.}
\end{exo}
\begin{exo}{}
A block rests on an incline. The static friction on it points
\choices{down the incline}{into the surface}{nowhere: it is zero}{up the incline}
\answer{D}{It keeps the block from sliding down.}
\end{exo}

%%DOC Q2.8 | Spring Forces
\begin{exo}{}
A spring with $k=100\Nm$ is stretched by $0.3\m$. What force does it exert?
\choices{$30\N$}{$300\N$}{$3\N$}{$333\N$}
\answer{A}{$100(0.3)=30\N$.}
\end{exo}
\begin{exo}{}
A force of $8\N$ stretches a spring by $4\,\mathrm{cm}$. What is its spring constant?
\choices{$2\Nm$}{$0.5\Nm$}{$200\Nm$}{$32\Nm$}
\answer{C}{$8/0.04=200\Nm$.}
\end{exo}
\begin{exo}{}
A $1\kg$ mass hangs at rest and stretches a spring by $0.02\m$. What is $k$?
\choices{$50\Nm$}{$500\Nm$}{$20\Nm$}{$5000\Nm$}
\answer{B}{$k=mg/\Delta x=10/0.02=500\Nm$.}
\end{exo}
\begin{exo}{}
A compressed spring pushes on the object at its end
\choices{toward the compression}{outward, back toward its natural length}{always upward}{not at all}
\answer{B}{Restoring force.}
\end{exo}
\begin{exo}{}
On a force vs.\ stretch graph, a spring's line passes through $(0.1\m,\,5\N)$. What is $k$?
\choices{$0.5\Nm$}{$5\Nm$}{$50\Nm$}{$500\Nm$}
\answer{C}{Slope $=5/0.1=50\Nm$.}
\end{exo}
\begin{exo}{}
The stretch of a spring triples. The spring force
\choices{stays the same}{becomes one third}{becomes nine times larger}{triples}
\answer{D}{Proportional to the stretch.}
\end{exo}

%%DOC Q2.9 | Circular Motion
\begin{exo}{}
An object moves in a circle at constant speed. Is it accelerating?
\choices{No, its speed is constant.}{Yes, toward the center.}{Yes, away from the center.}{Yes, along the tangent.}
\answer{B}{The direction of the velocity changes.}
\end{exo}
\begin{exo}{}
What is the centripetal acceleration for $v=6\ms$ and $r=3\m$?
\choices{$2\mss$}{$18\mss$}{$12\mss$}{$6\mss$}
\answer{C}{$36/3=12\mss$.}
\end{exo}
\begin{exo}{}
A $2\kg$ object moves in a circle of radius $2\m$ at $4\ms$. What net force acts on it?
\choices{$16\N$}{$8\N$}{$4\N$}{$32\N$}
\answer{A}{$2(16)/2=16\N$.}
\end{exo}
\begin{exo}{}
What provides the centripetal force that keeps the Moon in orbit?
\choices{Friction}{A normal force}{Tension}{Earth's gravity}
\answer{D}{Gravity points toward Earth, the center of the orbit.}
\end{exo}
\begin{exo}{}
The speed is halved at the same radius. The centripetal force becomes
\choices{half as large}{one quarter as large}{twice as large}{the same}
\answer{B}{$(1/2)^2=1/4$.}
\end{exo}
\begin{exo}{}
An object goes around a circle of radius $2\m$ in $4\s$. What is its speed?
\choices{About $3.14\ms$}{About $6.28\ms$}{About $1.57\ms$}{$8\ms$}
\answer{A}{$2\pi(2)/4=\pi\approx3.14\ms$.}
\end{exo}

%%DOC QTEST | Unit 2 test \textbullet\ Force and Translational Dynamics
\begin{exo}{}
A $3\kg$ mass is at $x=0$ and a $1\kg$ mass is at $x=8\m$. Where is their center of mass?
\choices{$x=4\m$}{$x=2\m$}{$x=6\m$}{$x=1\m$}
\answer{B}{$\frac{0+8}{4}=2\m$.}
\end{exo}
\begin{exo}{}
Which forces act on a ball after it leaves the thrower's hand (no air resistance)?
\choices{Gravity only}{Gravity and the force of the throw}{The force of the throw only}{None}
\answer{A}{The hand no longer touches the ball.}
\end{exo}
\begin{exo}{}
Earth pulls you down with $600\N$. You pull Earth with
\choices{$0$}{$600\N$, upward}{$6\N$}{$600\N$, downward}
\answer{B}{Third law.}
\end{exo}
\begin{exo}{}
Which must be true of an object in equilibrium?
\choices{It is at rest.}{No forces act on it.}{The net force on it is zero.}{It is moving.}
\answer{C}{It may be at rest or moving at constant velocity, with several forces acting.}
\end{exo}
\begin{exo}{}
Forces of $50\N$ east and $20\N$ west act on a $10\kg$ object. What is its acceleration?
\choices{$7\mss$ east}{$3\mss$ east}{$5\mss$ east}{$3\mss$ west}
\answer{B}{Net $30\N$ east; $30/10=3\mss$.}
\end{exo}
\begin{exo}{}
An elevator moving up slows down with an acceleration of $2\mss$ (downward). What does a scale read for a $50\kg$ person?
\choices{$400\N$}{$500\N$}{$600\N$}{$100\N$}
\answer{A}{The acceleration is downward: $50(10-2)=400\N$.}
\end{exo}
\begin{exo}{}
A $4\kg$ block on a frictionless table is tied to a hanging $1\kg$ block. What is the acceleration?
\choices{$2.5\mss$}{$10\mss$}{$8\mss$}{$2\mss$}
\answer{D}{$\frac{10}{4+1}=2\mss$.}
\end{exo}
\begin{exo}{}
The distance between two masses is halved. The gravitational force becomes
\choices{twice as large}{half as large}{four times as large}{one quarter as large}
\answer{C}{$1/(1/2)^2=4$.}
\end{exo}
\begin{exo}{}
A $30\kg$ box is pulled at constant velocity on a level floor with $\mu_k=0.2$. What horizontal pull is needed?
\choices{$60\N$}{$300\N$}{$6\N$}{$0$}
\answer{A}{Constant velocity: pull $=f_k=0.2(300)=60\N$.}
\end{exo}
\begin{exo}{}
A $5\kg$ box rests on a level floor with $\mu_s=0.6$. What is the maximum static friction?
\choices{$3\N$}{$50\N$}{$30\N$}{$8.3\N$}
\answer{C}{$0.6(50)=30\N$.}
\end{exo}
\begin{exo}{}
A spring with $k=250\Nm$ is stretched by $0.04\m$. What force does it exert?
\choices{$10\N$}{$6250\N$}{$1\N$}{$100\N$}
\answer{A}{$250(0.04)=10\N$.}
\end{exo}
\begin{exo}{}
A spring scale reads $15\N$ when its spring is stretched $0.05\m$. What is $k$?
\choices{$0.75\Nm$}{$3\Nm$}{$300\Nm$}{$30\Nm$}
\answer{C}{$15/0.05=300\Nm$.}
\end{exo}
\begin{exo}{}
A $0.5\kg$ ball on a string moves in a horizontal circle of radius $2\m$ at $4\ms$ (ignore gravity). What is the tension?
\choices{$1\N$}{$8\N$}{$16\N$}{$4\N$}
\answer{D}{$0.5(16)/2=4\N$.}
\end{exo}
\begin{exo}{}
A car's maximum speed on a flat curve is $v=\sqrt{\mu_sgr}$. If the car's mass doubles, its maximum speed
\choices{doubles}{is halved}{stays the same}{quadruples}
\answer{C}{The mass cancels: friction and the required force both double.}
\end{exo}
\begin{exo}{}
A box slides down a frictionless $30^\circ$ incline. What is its acceleration?
\choices{$10\mss$}{$5\mss$}{$8.7\mss$}{$0$}
\answer{B}{$g\sin30^\circ=5\mss$.}
\end{exo}
